\section{Exclusion of the obstruction}\label{sec:conclusion}

We apply Theorem~\ref{thm:exterior-count} to the exterior planes in
Proposition~\ref{prop:rig-obstruction}.  Four small exterior coordinates
give two small signed products, whose range in the trace plane has area
of order $r^4$.  The integral along a diagonal orbit contributes
logarithmic factors; we will show that they are summable at this scale.

\subsection{The exterior vector integral}

Fix a partition $b_1=ij|kl$ of $\{1,2,3,4\}$, and denote the other two
partitions by $b_2,b_3$.  For $m\in\Z\setminus\{0\}$, $L>1$, and
$0<r<1$, define $N_{m,L,r}(\Lambda)$ to be the number of
$v\in\bigwedge^2\Lambda$ satisfying
\begin{equation}\label{eq:end-vector-box}
 Q(v)=m,\qquad |v_{ij}|,|v_{kl}|<L,\qquad
 |v_{ab}|<r\quad(ab\ne ij,kl).
\end{equation}
This count is finite because $\bigwedge^2\Lambda$ is a lattice in
$\bigwedge^2\R^4$ and the specified box is bounded.  To see that it is
Borel, choose a local continuous determinant-one basis around a lattice
with $k$ witnesses and keep their distinct integer exterior labels fixed.
They remain distinct, and their $Q$-values remain $m$ because
$Q((\bigwedge^2g)w)=\det(g)Q(w)$.  Their coordinates vary continuously,
so the strict inequalities persist in a neighborhood.  Thus
$\{N_{m,L,r}\ge k\}$ is open for every integer $k\ge1$.
Recall that the three real embeddings of $F$ are indexed by
$\mathcal B$.  For a nonzero $z\in F$, set
\begin{equation}\label{eq:end-logarithms}
 \mathcal L_{L,r}(z)
  =\log_+\frac{L^2}{|z_{b_1}|}\,
   \log_+\frac{r^2}{|z_{b_2}|}\,
   \log_+\frac{r^2}{|z_{b_3}|}.
\end{equation}
All three conjugates of a nonzero element of the field $F$ are nonzero,
so the expression is defined.

In Equation~\eqref{eq:ext-multiplicity}, $M(z)$ sums over
the fixed maximal-order ideal representatives $I$ and the $U^+$-orbits
of nonzero raw $w\in\bigwedge^2_{\Z}I$ with $\Psi(w)/\Nm I=z$.  Each
orbit has weight $\prod_pP_p(I,w)$, the average membership of $w$ in
$\bigwedge^2_{\Z}L'$ over the local labels
$L'\in\mathscr L_{\OO}(I)$.

\begin{lemma}[Exterior unfolding]\label{lem:end-unfold}
For every primitive totally real quartic field $K$, order $\OO$, and
ordering of its real embeddings, let $d_U=[U:U^+]$.  With the preceding
choices of $b_1,m,L,r$, use the same fixed system of maximal-order ideal
representatives in the packet model and in the definition of $M$.  Then
\begin{equation}\label{eq:end-unfold}
 \int_{X_4}N_{m,L,r}(\Lambda)\,d\mu_{\OO,\sigma}(\Lambda)
 =\frac{4}{d_U\sqrt D\,\kappa}
       \sum_{t\in\{m,-m\}}
       \sum_{\substack{z\in F\\\Tr z=t}}
          M(z)\mathcal L_{L,r}(z).
\end{equation}
The index satisfies $2\le d_U\le16$.
\end{lemma}

\begin{proof}
We use the packet integral in Equation~\eqref{eq:pkt-integration}.
For its maximal-order ideal representative $I$ and local label
$L'\in\mathscr L_{\OO}(I)$, the normalized lattice at the sign vector
$s=(s_1,\ldots,s_4)\in\{\pm1\}^4$ and the logarithmic parameter
$(u_1,\ldots,u_4)\in\R^4$ with $u_1+\cdots+u_4=0$ is
\[
 (X\Nm I)^{-1/4}\sigma(L')
       \diag(s_1e^{u_1},\ldots,s_4e^{u_4}).
\]
For a raw vector $w\in\bigwedge^2 I$, its normalized exterior coordinates
are
\[
 y_{ab}=(X\Nm I)^{-1/2}s_as_b e^{u_a+u_b}w_{ab}.
\]
If $z=\Psi(w)/\Nm I$ and $\epsilon(s)=\prod_a s_a$, then
Equation~\eqref{eq:ext-signed-products} gives
\begin{equation}\label{eq:end-signed-products}
 (y_{12}y_{34},-y_{13}y_{24},y_{14}y_{23})
       =\epsilon(s)z,\qquad
 Q(y)=\epsilon(s)\Tr z.
\end{equation}
There are eight sign vectors for each value of $\epsilon(s)$.

The kernel of the signature map on $U$ is $U^+$, so
$d_U=|\operatorname{sgn}U|\le16$; the signature of $-1$ gives
$d_U\ge2$.  After averaging over $L'$, the vector-count integrand in
Equation~\eqref{eq:pkt-integration} is invariant under
$v\mapsto\sigma(e)v$ for $e\in U$, since $L'\mapsto eL'$ permutes the
labels and transports their exterior lattices.  Replacing
the fundamental domain for $U$ in $T_\infty$ by one for $U^+$ therefore
multiplies its integral by $d_U$, which we compensate by division by
$d_U$.  The latter domain consists of all sign components, with a
fundamental domain for the positive unit logarithms in each component.

For a given raw $w$, its average membership over the local labels is
\[
 \frac1{N_{\OO}}\sum_{L'\in\mathscr L_{\OO}(I)}
       \mathbf 1_{\{w\in\bigwedge^2L'\}}
     =\prod_p P_p(I,w),
\]
the weight in Equation~\eqref{eq:ext-multiplicity}.  This weight is constant on its
$U^+$-orbit.  That action is free: every embedding coordinate of the
nonzero raw vector $w\in\bigwedge^2_{\Z}I\subset\bigwedge^2_{\Q}K$ is
nonzero by Lemma~\ref{lem:ext-products}.  A stabilizing positive unit
$e$ would therefore satisfy $\sigma_a(e)\sigma_b(e)=1$ for every pair
$ab$, forcing $e=1$.  Summing the translates of each raw vector thus unfolds the
positive-unit fundamental domain to all of the three-dimensional space
$u_4=-u_1-u_2-u_3$, with multiplicity one.  The packet normalization
remaining after the uniform label average is
$h_K\vol(\mathscr F_U)=\sqrt D\,\kappa$ by
Equation~\eqref{eq:pkt-unit-volume}.  Hence this unfolding has coefficient
$1/(d_U\sqrt D\,\kappa)$ before the sign sum and the integral in $u$.
All the sums and integrals are nonnegative, so Tonelli's theorem applies.

Choose one coordinate from each opposite pair.  The logarithms of the
absolute values of these coordinates are affine functions of
$(u_1,u_2,u_3)$.  For the choices $12,13,14$, their linear parts are
\[
 u_1+u_2,\qquad u_1+u_3,\qquad -u_2-u_3.
\]
Their matrix has determinant of absolute value two.  Choosing the
opposite member of a pair negates its row, and ordering the partitions
only permutes rows, so the same absolute determinant holds for every
choice.  If the two entries in an opposite pair have product of absolute
value $|z_b|$ and both must be less than $h>0$ in absolute value, the
logarithm of the chosen entry lies in an interval of length
$\log_+(h^2/|z_b|)$.  Therefore, for a raw vector with the required
trace after the sign change, its integral over $(u_1,u_2,u_3)$ is
$\frac12\mathcal L_{L,r}(z)$.

The $U^+$-orbit sum of the weights, including the sum over $I$, is exactly
$M(z)$ by Equation~\eqref{eq:ext-multiplicity}.  Equation
\eqref{eq:end-signed-products} allows raw trace $m$ for the eight positive
determinant signs and raw trace $-m$ for the other eight.  Multiplying
eight by the Jacobian factor $1/2$ gives the factor $4$ in
Equation~\eqref{eq:end-unfold}.  The unit $-1$, although it acts trivially
on exterior vectors, has already been accounted for in the index $d_U$.
\end{proof}

\begin{proposition}[Small exterior boxes]\label{prop:end-tube}
For every fixed $m\in\Z\setminus\{0\}$ and $L>1$, there are constants
$r_0>0$ and $C_{m,L}>0$ with the following property.  For every sequence
of the fields, orders, and embedding orderings in Theorem~\ref{thm:main}
with $X_i\to\infty$, and every fixed $0<r<r_0$,
\begin{equation}\label{eq:end-tube}
 \limsup_{i\to\infty}
   \int_{X_4}N_{m,L,r}(\Lambda)\,
           d\mu_{\OO_i,\sigma_i}(\Lambda)
       \le C_{m,L}r^4.
\end{equation}
The constants are independent of the sequence, of $r$, and of the choice
of partition $b_1$.
\end{proposition}

\begin{proof}
Take $r_0<1$ so small that $2r_0^2\le |m|/2$.  For each
$t\in\{m,-m\}$, fix the coordinate rectangle
\[
 \Omega_t=\{z\in\mathcal H_t:|z_{b_2}|,|z_{b_3}|\le1\}.
\]
Now let $0<r<r_0$.  In the trace-$t$ sum in
Equation~\eqref{eq:end-unfold}, only points with
$|z_{b_2}|,|z_{b_3}|<r^2$ contribute, and these points lie in $\Omega_t$.
Since their three coordinates have sum $t$,
\[
 |z_{b_1}|\ge |m|-2r^2\ge |m|/2.
\]
Thus the first logarithm in Equation~\eqref{eq:end-logarithms} is bounded
in terms of $m,L$ on the contributing points, independently of $r$.
For every $z\in\Omega_t$ with $M(z)>0$, Lemma~\ref{lem:ext-support} gives
$|z_b|\gg_{\Omega_t}X^{-84}$ for each $b$.  Hence the last two logarithms
in Equation~\eqref{eq:end-logarithms} are $O(\log X)$, uniformly for
$0<r<1$.

We first remove the points in the two strips
\begin{equation}\label{eq:end-strips}
 |z_{b_2}|\le4X^{-\tau}
       \quad\text{or}\quad
 |z_{b_3}|\le4X^{-\tau}.
\end{equation}
For large $X$, each strip within $\Omega_t$ is covered by a rectangle
of width $8X^{-\tau}$ and length $2$ in the coordinates
$z_{b_2},z_{b_3}$, and these rectangles lie in $\Omega_t$.  Both sides
meet the minimum length in
Theorem~\ref{thm:exterior-count} for large $X$.  That theorem bounds the
sum of $M(z)$ over their union, divided by $\sqrt D\,\kappa$, by
$O(X^{-\tau})$.  The logarithmic bounds just proved show that their
contribution to Equation~\eqref{eq:end-unfold} is
\[
 O_{m,L}\bigl(X^{-\tau}(\log X)^2\bigr)=o(1).
\]
The use of the rectangle estimate here controls the weighted mass
directly; it does not require a pointwise bound for $M(z)$.

For the remaining points, separate the signs of $z_{b_2},z_{b_3}$ and
use the dyadic intervals
\[
 r^2 2^{-a-1}<|z_{b_2}|\le r^2 2^{-a},\qquad
 r^2 2^{-b-1}<|z_{b_3}|\le r^2 2^{-b},
 \qquad a,b\ge0.
\]
Each signed dyadic rectangle lies in $\Omega_t$.
Each interval has width half its upper endpoint.  If it meets
$|z_{b_j}|>4X^{-\tau}$, that width exceeds $2X^{-\tau}$.  Hence every
dyadic rectangle meeting the remaining region satisfies the minimum
side lengths of Theorem~\ref{thm:exterior-count}; any boundary rectangle
may be included in full.  Its area is $O(r^4 2^{-a-b})$, while the last
two factors of Equation~\eqref{eq:end-logarithms} are at most
$(a+1)\log2$ and $(b+1)\log2$.  Theorem~\ref{thm:exterior-count} gives
\[
 \frac1{\sqrt D\,\kappa}
   \sum_{z\text{ in that rectangle}} M(z)
       \ll_{m,L} r^4 2^{-a-b}.
\]
The sum over the four choices of signs and over all $a,b$ is therefore
at most
\[
 C_{m,L}r^4
       \sum_{a,b\ge0}(a+1)(b+1)2^{-a-b}
 \ll_{m,L}r^4.
\]
All constants are independent of $r$: the rectangles $\Omega_t$ and the
uniform rectangle constant were fixed before choosing $r$.  The two
strips tend to zero for each fixed $r$, and $4/d_U\le2$ in
Equation~\eqref{eq:end-unfold}.  These observations prove
Equation~\eqref{eq:end-tube} for both traces and all three partitions.
\end{proof}

\subsection{Identification of every probability limit}

\begin{proof}[Proof of Theorem~\ref{thm:main}]
For the sequence in that theorem,
$X_i=\sqrt{|\Disc(\OO_i)|}$ tends to infinity.  By
Proposition~\ref{prop:ord-tightness}, the packet probabilities are tight,
and every subsequence has a further weakly convergent subsequence with
an $A_4$-invariant probability limit $\mu$.
Proposition~\ref{prop:ord-balls} gives the local ball bound required by
Corollary~\ref{cor:rig-positive-components}.  Hence almost every component in
the $A_4$-ergodic decomposition of $\mu$ has positive entropy for a
diagonal flow.  Proposition~\ref{prop:rig-obstruction} reduces the
identification of these components to the exclusion of its three
exterior-plane events.

Fix a partition $b_1$, a nonzero integer $m$, and $L>1$.  The event
\[
 E_{m,L,r}=\{\Lambda\in X_4:N_{m,L,r}(\Lambda)\ge1\}
\]
is open by the observation following Equation~\eqref{eq:end-vector-box}.

The open-set inequality for weak convergence and
Proposition~\ref{prop:end-tube} now give, along the convergent subsequence,
\[
 \mu(E_{m,L,r})
  \le\liminf_i\mu_{\OO_i,\sigma_i}(E_{m,L,r})
  \le\limsup_i\int N_{m,L,r}\,d\mu_{\OO_i,\sigma_i}
  \le C_{m,L}r^4
\]
for every sufficiently small fixed $r>0$.  In particular,
\[
 \mu\left(\bigcap_{n\ge2}E_{m,L,1/n}\right)
 \le \inf_{\substack{n\ge2\\1/n<r_0}} C_{m,L}n^{-4}=0.
\]
Every vector in one of the obstruction planes with $Q\ne0$ has a
nonzero integer $Q$-value by Equation~\eqref{eq:rig-Q-integral}, and its
two entries are less than some integer $L\ge2$ in absolute value.
An exact witness belongs to $E_{m,L,1/n}$ for every $n\ge2$.  Thus the
Borel event in Equation~\eqref{eq:rig-events} satisfies
\[
 \mathcal E_{b_1}
 \subset
 \bigcup_{m\in\Z\setminus\{0\}}\ \bigcup_{\substack{L\in\Z\\L\ge2}}
       \bigcap_{n\ge2}E_{m,L,1/n}.
\]
The right side is a countable union of null sets, so
$\mu(\mathcal E_{b_1})=0$.  The same argument applies to
all three partitions.  Integrating these null events over the ergodic
decomposition shows that the proper alternative in
Proposition~\ref{prop:rig-obstruction} has zero total component weight.
Almost every component is therefore $m_4$, and hence $\mu=m_4$.

Every subsequential probability limit is now identified as $m_4$.
Tightness implies that the entire sequence converges weakly to $m_4$,
which gives the asserted convergence against every $f\in C_c(X_4)$.
The same tightness statement supplies, for each $\epsilon>0$, one
compact set carrying at least $1-\epsilon$ of every sufficiently late
packet probability.  This is the required no-escape assertion.
\end{proof}
